% ------------------------------------------------------------------
%  4.2.5 Especificaciones Implementos a proveer (Hardware y Software)
%  Fragmento del Subdocumento 4.2, traido desde contenido.tex con
%  \parte{partes/04_c_implementos}. Las rutas son relativas a la carpeta del
%  subdocumento, nunca a la de main.tex.
%
%  Resumen y análisis; el detalle fila por fila vive en el Anexo 4.A
%  (Formulario T-11, partes/15_anexo_t11).
% ------------------------------------------------------------------

\section{Especificaciones Implementos a proveer (Hardware y Software)}
\label{sec:4-2-6-implementos}

Esta sección resume y analiza el equipamiento y el software que la solución requiere en las instalaciones del CLIENTE, en el terreno y en la plataforma de nube. El detalle de cada elemento, con su producto ofertado, su ubicación, su cantidad y su justificación, se entrega en el Anexo 4.A (Formulario T-11), que acompaña a esta parte como anexo. El equipamiento físico lo adquiere el CLIENTE, y el adjudicatario lo especifica, instala, integra y mantiene durante el contrato (Art. 14.2).

\subsection{Síntesis del equipamiento por familia}

La Tabla~\ref{tab:t26} agrupa el equipamiento en ocho familias, con sus cantidades y los sitios donde se instala.

\begin{tablalafrox}{Síntesis del equipamiento a proveer por familia}{tab:t26}%
  {F{0.22} Y F{0.26}}%
  {\cab{Familia} & \cab{Equipamiento y cantidad} & \cab{Sitios}}
  Cómputo & 3 servidores del clúster, 1 servidor de borde y 3 mini-PC industriales & Talca, Concepción y cross-docking \\
  Almacenamiento de respaldo & 1 NAS WORM cifrado, D-05 & Talca \\
  Red y seguridad & 6 firewalls, 6 switches de sitio, 1 switch de gestión y 9 puntos de acceso Wi-Fi 6E & Los cinco sitios con cómputo \\
  Enlaces & 2 enlaces de fibra, 2 enlaces LTE en los CD, 3 módems LTE de doble SIM con dos proveedores en los cross-docking y 3 kits Starlink & CD y cross-docking \\
  Terreno & 62 terminales EC55 más 20 \% de reserva, cerca de 200 terminales TC58e, 200 impresoras de cabina y 200 terminales de pago & Calle y camiones \\
  Bodega & 174 terminales de cámara, 6 impresoras de andén, 3 balanzas y 6 escáneres GS1 & CD y cross-docking \\
  Cadena de frío & 7 módulos de sensores con 28 puntos, 2 gateways IoT y 28 termógrafos & CD, cámaras y camiones \\
  Puestos y sala técnica & 8 estaciones de trabajo y la infraestructura de sala de Talca: UPS, generador, 2 unidades de clima de precisión, 4 gabinetes y 7 cámaras & Talca y Concepción \\
\end{tablalafrox}
\fuente{elaboración propia; detalle en el Anexo 4.A (Formulario T-11).}

Más de nueve de cada diez unidades están en el terreno y en la bodega, no en la sala técnica: son unos 860 dispositivos de trabajo frente a menos de 50 equipos de cómputo, red y sala. Esa proporción refleja el caso, donde la operación ocurre en la calle, en la cámara y en el andén, y explica por qué la gestión de dispositivos (N-13) es un componente con emplazamiento propio y no un accesorio.

\subsection{Criterios de selección}

Cada familia se especifica contra una condición del caso que el equipamiento de oficina no resiste. Los 174 terminales de cámara están certificados para trabajar a $-$30\,\textdegree{}C con guantes gruesos y batería de recambio en caliente, porque la preparación de congelados ocurre a $-$22\,\textdegree{}C (Caso, Cap. 4.5); su número cubre los 120 preparadores simultáneos de Talca con 20 \% de reserva por la rotación de 38 \%, más 30 en Concepción. Los terminales de terreno operan un turno completo sin señal, a una mano y bajo sol directo (RT-13.08). Los mini-PC de los cross-docking son de rango industrial, porque las naves no tienen climatización y el supuesto de diseño es de hasta 50\,\textdegree{}C bajo la techumbre en verano. Los 28 termógrafos cubren los 18 camiones propios y los 10 refrigerados de los transportistas que declara la Tabla~\ref{tab:t73} del caso, porque el instrumento viaja con la carga y su registro ampara la cadena de frío aun cuando el vehículo no sea del CLIENTE; están calibrados contra un patrón NIST en dos puntos, porque su registro es la prueba de la cadena de frío ante un cliente o ante la autoridad sanitaria.

En la sala técnica el criterio es la redundancia sin sobrecompra. Los tres servidores del clúster son idénticos, y el servidor de Concepción es del mismo modelo, de modo que la pérdida de cualquiera deja capacidad para todas las máquinas virtuales y Concepción puede asumir el WMS de Talca, y la cuarta unidad solo se agrega con el crecimiento de 3 veces, en el gabinete R04 ya reservado. Los firewalls de Talca van en par, y los de Concepción y de cada cross-docking son únicos, porque en esos sitios la contingencia ante la falla es la misma autonomía que cubre los cortes de enlace.

\subsection{Software y licenciamiento}

El software de base es de código abierto en todos los componentes que el adjudicatario opera: PostgreSQL, RabbitMQ, Keycloak, Proxmox VE, Ceph y Docker, además de la aplicación, escrita en Django y Python y, en el terreno, en Kotlin. Los únicos elementos con suscripción son los servicios de la plataforma de nube, descritos en la sección~\ref{sec:4-2-2-servicios-nube}, la gestión de dispositivos y los agentes de detección en endpoints. Esta composición mantiene la reversibilidad de la solución y evita que el CLIENTE quede atado a un licenciamiento propietario. El Anexo 4.A (Formulario T-11) detalla, además del equipamiento, cada servicio de nube y cada componente de software con su cantidad y su justificación.
